\section{Conclusion and Limitations}

Moving this controller onto a vehicle is a systems problem rather than a control
problem. The parts are a phone the driver already owns, a US\$250 single-board
computer, a traffic API, and a filter on the advisory, with nothing installed on
the road and no server inside the control loop. The compute device's median is
19.8\,ms and the link's 95th percentile is 116.19\,ms against a 200\,ms target.
Restricting the observation to what a feed returns changes throughput by
6\,veh/h against error bars of 44 and 61\,veh/h, while both controlled arms gain
more than 200\,veh/h over no control. What decided whether a drive produced data
was thermal, physical and perceptual rather than any of this.

\noindent\textbf{What the deployment does not establish.}
No advisory was acted on, so nothing here measures compliance. The corpus is
eight runs on one day, with one vehicle, one driver and one road. The frame
rotation left the detector returning almost nothing, so the safety layer's
camera inputs were processed offline. The wire hop is an order rather than a
duration, because its median and its conversion bound are the same size. The
feed reports when a response arrived and not when the conditions held, so
nothing bounds the observation's staleness. The first drive ended in a condition
nothing detects, and distance is a bracket because the two methods disagree by
3.5\,km.

\noindent\textbf{What the simulation does not establish.}
Every figure is against this port's own baseline and none is comparable to the
published table. The port differs from its source in right-of-way, exit geometry
and gradient, the fleet is not the published one, the episode does not reach a
steady state, and the jam factor the simulated observation reads is a stated
model rather than a measurement.

\section{Acknowledgements}
Opus 5 (1M) model was used for writing the code as well as the text of the paper.
All sections in the paper and code was co-worked on by AI and human authors, with human verification.